% ============================================================
% Resumen (ES) y Abstract (EN)
% ============================================================

\phantomsection
\addcontentsline{toc}{chapter}{Resumen}

\begin{flushleft}
{\Huge \textbf{\sffamily Resumen}}
\end{flushleft}
\vspace{0.5cm}

Este Trabajo de Fin de Grado presenta el diseño, implementación y validación de un sistema de vigilancia distribuida basado en una red de tres nodos radar cooperativos sobre plataformas embebidas de bajo coste. Se ha construido un prototipo físico funcional en el que cada nodo integra un microcontrolador ESP32 de doble núcleo, un sensor ultrasónico HC-SR04 montado sobre un motor paso a paso NEMA~17 y una estructura mecánica diseñada e impresa en 3D.

El problema central es la interferencia acústica cruzada (\textit{crosstalk}) que surge cuando varios transductores operan simultáneamente en el mismo espacio físico. Para resolverlo sin hardware analógico especializado, se diseña un protocolo de acceso al medio basado en TDMA dinámico con reutilización espacial: un servidor central calcula en tiempo real qué pares de nodos son geométricamente incompatibles y asigna ranuras temporales diferenciadas únicamente a aquellos cuyas zonas de emisión se solapan, permitiendo que el resto dispare en paralelo.

El firmware, implementado sobre FreeRTOS, dedica un núcleo exclusivamente a la pila de red Wi-Fi y el otro al control determinista del motor y a la medición del tiempo de vuelo mediante interrupciones de hardware. El servidor central, desarrollado en Python, orquesta el ciclo de supertrama, gestiona desconexiones y reincorporaciones de nodos de forma tolerante a fallos, y ejecuta un algoritmo de trilateración cooperativa para estimar la posición de los objetos detectados en zonas de solapamiento. Los resultados son visualizados en tiempo real a través de una interfaz gráfica desarrollada con Tkinter.

La validación experimental sobre el prototipo construido demuestra que el protocolo elimina por completo las lecturas espurias debidas al \textit{crosstalk} en todos los escenarios de prueba, que el sistema mantiene la cobertura del perímetro ante pérdidas y reconexiones de nodos sin intervención manual, y que las estimaciones de posición cooperativas son coherentes con la geometría del despliegue, confirmando la viabilidad del enfoque para redes de sensores acústicos de bajo coste.

\textbf{Palabras clave:} ESP32, radar ultrasónico, FreeRTOS, TDMA dinámico, trilateración,
vigilancia distribuida, redes de sensores inalámbricos.

\clearpage

% ============================================================

\phantomsection
\addcontentsline{toc}{chapter}{Abstract}

\begin{flushleft}
{\Huge \textbf{\sffamily Abstract}}
\end{flushleft}
\vspace{0.5cm}

This Final Degree Project presents the design, implementation, and validation of a distributed surveillance system based on a network of three cooperative radar nodes built on low-cost embedded platforms. A fully functional physical prototype has been constructed in which each node integrates a dual-core ESP32 microcontroller, an HC-SR04 ultrasonic sensor mounted on a NEMA~17 stepper motor, and a 3D-printed mechanical structure.

The core challenge is the acoustic cross-interference (\textit{crosstalk}) that arises when multiple ultrasonic transducers operate simultaneously in the same physical space. To solve it without specialised analogue hardware, a medium-access protocol based on dynamic TDMA with spatial reuse is designed: a central server computes in real time which pairs of nodes are geometrically incompatible and assigns separate time slots only to those whose emission cones overlap, while the remaining nodes fire in parallel.

The firmware, implemented on FreeRTOS, dedicates one core exclusively to the Wi-Fi stack and the other to deterministic motor control and hardware-interrupt-based time-of-flight measurement. The central server, developed in Python, orchestrates the superframe cycle, handles node disconnections and reconnections in a fault-tolerant manner, and runs a cooperative trilateration algorithm to estimate the position of objects detected in overlap zones. Results are visualised in real time through a Tkinter graphical interface.

Experimental validation on the built prototype shows that the protocol completely eliminates spurious readings caused by \textit{crosstalk} across all test scenarios, that the system maintains perimeter coverage under node loss and reconnection without manual intervention, and that cooperative position estimates are consistent with the deployment geometry, confirming the viability of the proposed approach for low-cost acoustic sensor networks.

\textbf{Keywords:} ESP32, ultrasonic radar, FreeRTOS, dynamic TDMA, trilateration,
distributed surveillance, wireless sensor networks.

\clearpage
